% Methods and data. Numbers: ../FINDINGS.md.
\section{Data}\label{sec:data}

\paragraph{Citation graph.} Citations come from the parsed OpenCitations snapshot of January 2024 released through
the Illinois Data Bank \citep{arjomandi2024dataset}, built from the OpenCitations Index \citep{peroni2020opencitations,heibi2019coci}.
It holds 1,962,840,983 citation links between 107,030,533 bibliographic resources. Identifiers (DOI, PMID) and journals
were taken from the OpenCitations Meta dump \citep{massari2024meta,ocmeta2024dump}: 26,582,951 graph nodes carry a PMID
and 91,379,362 a DOI. Counting records and citing links by publication month shows that the snapshot is essentially
complete up to September 2023: 240,000--300,000 records and 8.7--10.3 million citing links per month from January to
September 2023, then 58,434 records and 2.4 million links in October and fewer afterwards.

\paragraph{Retractions.} Retraction labels come from the Retraction Watch database distributed by Crossref
\citep{rittman2025rw}. One copy was frozen for the whole study (repository commit of 30 October 2026, 73,519 records).
Only records whose nature is \emph{Retraction} are used; expressions of concern and corrections are ignored. Records
were grouped into papers by DOI, or by PMID when no DOI is given, keeping the earliest retraction notice. This gives
62,735 retracted papers. A paper is labelled \emph{paper mill} when any of its Retraction Watch reasons contains ``Paper
Mill'' (11,716 papers).

\paragraph{PubMed.} The 2026 PubMed baseline (39,994,988 records) supplies PMID--DOI links, the PubMed retraction flags
and grant agencies. A paper is \emph{NIH-funded} when one of its grants lists an agency containing ``NIH''
(3,341,220 records). A graph node is \emph{biomedical} when it has a PMID in OpenCitations Meta or its DOI appears in
PubMed (31,464,435 nodes).

\paragraph{Matching.} Retraction Watch papers were matched to graph nodes by DOI and by PMID. Of the 62,735 retracted
papers, 52,254 (83.3\%) match at least one node. Among Retraction Watch papers with a PMID, 88.8\% are also flagged as
retracted in PubMed; Retraction Watch is used as the single label source and PubMed as a check. In a random sample of 300 DOI-matched papers,
the Retraction Watch and OpenCitations titles agreed for all 292 papers that had a title in OpenCitations Meta
(differences were retraction prefixes such as ``Notice of Retraction'' or truncation).

\section{Methods}\label{sec:methods}

\subsection{Design}
The test is prospective in the sense that matters for a screening tool: all predictors use only information available
at a fixed time point $T_0$, and outcomes are retractions issued after $T_0$. The main time point is
$T_0 = $ 31 January 2024. Only 116,123 citing links in the graph (0.006\%) carry a publication date after January 2024,
so post-$T_0$ information can enter the predictors only negligibly. Because the graph is complete only to about September 2023, every analysis was
repeated with $T_0' = $ 30 September 2023. The design, scores, metrics and decision rule were written into a
specification file whose SHA-256 hash was recorded in the lab notebook before any score was computed, and the analysis
code checks the hashes of the frozen score files before reading any outcome. The specification was not deposited with a
third-party registry before the outcomes were read, so the design is pre-specified rather than formally
pre-registered; analyses added later are labelled post hoc.

\paragraph{Candidates and outcome.} Candidates are biomedical nodes published 2015--2023 that have a journal, at least
one reference in the graph, and no retraction notice dated on or before $T_0$ (10,495,850 papers). A candidate is
\emph{positive} if its first retraction notice falls between $T_0$ and 30 October 2026. Let $R_0$ denote the 44,401
graph nodes retracted on or before $T_0$. $R_0$ is defined by the date of the retraction notice; the Retraction
Watch export does not record when a notice entered the database, so a few notices dated before $T_0$ may have been
recorded only later and would not have been visible to a screen run on that date.

\paragraph{Scores.} For a candidate $p$ with $d(p)$ references in the graph:
$S_1(p)$ is the number of references of $p$ that are in $R_0$;
$S_2(p) = S_1(p)/d(p)$;
$S_3(p)$ is the number of citations $p$ received from papers in $R_0$; and
$S_4(p)$ is the number of references of $p$ to paper-mill retractions in $R_0$.
The journal baseline is a smoothed retraction rate,
\[
B(p) = \frac{r_j + m\,g}{n_j + m},
\]
where $j$ is the journal of $p$, $n_j$ is the number of biomedical items in $j$ published from 2010 to 2023,
$r_j$ is the number of those in $R_0$, $g$ is the overall rate $\sum_j r_j / \sum_j n_j$, and $m = 100$.

\subsection{Analyses}
\paragraph{Association.} Positive rates with Wilson intervals \citep{wilson1927} are reported by $S_1$ level, and risk
ratios for $S_1 \ge 1$, $S_3 \ge 1$ and $S_4 \ge 1$.

\paragraph{Within journal (primary test).} Retractions cluster by journal and publication year, for example in
mass retractions of special issues. To separate the neighbourhood signal from these clusters, the Mantel--Haenszel odds
ratio \citep{mantel1959} of exposure ($S_1 \ge 1$) on outcome was computed with strata defined by journal $\times$
publication year, with the Robins--Breslow--Greenland confidence interval \citep{robins1986estimator}. The
pre-specified decision rule was that the signal holds within journal if the 95\% interval lies above 1.

\paragraph{Screening value.} Candidates were ranked by (i) $S_1$, with ties broken by $S_2$ and then at random;
(ii) the journal baseline $B$; and (iii) a logistic model with $\operatorname{logit} B$, $\log(1+S_1)$,
$\log(1+S_3)$, an indicator of $S_1 \ge 1$ and indicators of age at the time point. The model was fitted only on an
earlier window (predictors at 31 December 2021; outcomes up to $T_0$; candidates published 2013--2021) and applied
unchanged. Precision at $k$, lift (precision divided by the base rate), recall and average precision
\citep{saito2015prc} are reported with 95\% intervals from 1,000 Poisson bootstrap resamples of the candidates.

\paragraph{Sensitivity analyses.} The analyses were repeated (a) for candidates published 2015--2022, (b) without
nodes whose DOI maps to more than one graph node, and (c) with $T_0'$.

\paragraph{Supporting analyses.} Three descriptive analyses complement the test. (1) \emph{Community concentration}:
each retracted biomedical paper with a journal was matched to one non-retracted paper from the same journal and
publication year, chosen by a fixed hash of its identifier; communities were found with the Leiden algorithm \citep{traag2019leiden} (modularity, two iterations, seeds 1
and 2) on the subgraph formed by
these papers and their neighbours, and the concentration of retracted papers over communities was compared with that
of the controls under a permutation null that swaps labels within pairs. (2) \emph{NIH exposure}: the share of
biomedical papers published 2000--2023 that cite a paper in a year after its retraction, for NIH-funded and other
papers, with a Mantel--Haenszel odds ratio stratified by journal and year. (3) \emph{Citations around retraction}:
citations per paper by year relative to the retraction year, for papers retracted 2000--2021 and matched controls.

\paragraph{Implementation.} All joins and aggregations run in DuckDB \citep{raasveldt2019duckdb} on one 8-core, 64 GB
cloud machine. Every run writes its configuration, the hash of its script, its counts and checksums; per-paper scores
are never released. The retracted papers, the citation links that touch them and the community memberships are
released as a dataset \citep{arjomandi2026retracteddata}.
